\section{LLM, Context Window וזיכרון}

מודל השפה הגדול (\eng{LLM}) הוא ליבת הסוכן, אך חשוב להבין את מגבלותיו:
למודל אין זיכרון פנימי מתמשך בין הרצות, וכל המידע שהוא ``רואה'' מוגבל
לגודל חלון ההקשר (\eng{Context Window}) \cite{brown2020gpt3}.

\subsection{מהו LLM?}

\eng{LLM} (\eng{Large Language Model}) הוא רשת נוירונים מאומנת על כמויות
עצומות של טקסט, המסוגלת לחזות המשך סביר לרצף מילים. בפועל, המודל מבצע
משימות כמו סיכום, תרגום, כתיבת קוד ותכנון --- אך תמיד בתוך מסגרת
הטוקנים שהוזנו אליו בבקשה הנוכחית.

\subsection{Context Window --- זיכרון העבודה}

חלון ההקשר הוא ``שולחן העבודה'' של הסוכן בזמן שיחה. בו מתנהלים:
\begin{itemize}[rightmargin=2em]
  \item \eng{System Prompt} --- ההנחיות הכלליות לסוכן.
  \item היסטוריית השיחה --- שאלות ותשובות קודמות.
  \item תוצאות כלים --- פלט מחיפוש, הרצת קוד, שליפת \eng{RAG}.
  \item תוכן \eng{Skills} שנבחרו לבקשה הנוכחית.
\end{itemize}

כאשר החלון מתמלא, יש לקצץ, לסכם או להעביר מידע לאחסון חיצוני. זו אחת
הסיבות לחשיבות \eng{RAG} ולניהול זיכרון ארוך טווח.

\begin{notebox}
לפי ההרצאה: ``אין לו זיכרון בעצמו --- הכל נמצא אצלנו במחשב וכל פעם
אני טוען את ה-\eng{Context Window} מחדש עם כל מה שהיה עד עכשיו.''
\end{notebox}

\subsection{Memory --- סוגי זיכרון}

נהוג להבחין בין:
\begin{description}[rightmargin=2em, style=nextline]
  \item[זיכרון קצר טווח] תוכן חלון ההקשר הנוכחי.
  \item[זיכרון ארוך טווח] מאגר חיצוני --- קבצים, מסד נתונים, \eng{vector store}.
  \item[זיכרון פרוצדורלי] \eng{Skills}, תבניות \eng{CLS} וכללים קבועים.
\end{description}

\subsection{הרכבת פרומפט מלא}

לפני כל קריאה ל-\eng{LLM}, המערכת מרכיבה חבילה:
\[
  \text{Input} = \text{System} + \text{Skills} + \text{RAG Docs} + \text{User Prompt} + \text{Tool Outputs}
\]
המודל מייצר תשובה; התשובה מתווספת לחלון ההקשר; המחזור חוזר.

\subsection{אסטרטגיות לניהול הקשר}

\begin{enumerate}[rightmargin=2em]
  \item \textbf{סיכום תקופתי} --- סוכן מסכם שיחה ארוכה ומחליף היסטוריה בסיכום.
  \item \textbf{שליפה סלקטיבית} --- רק מסמכים רלוונטיים מ-\eng{RAG} נכנסים לחלון.
  \item \textbf{חלוקה למשימות} --- אורכוסטרטור מעביר תת-משימות לסוכנים עם הקשר מצומצם.
\end{enumerate}

\begin{insightbox}
\eng{Chain-of-Thought} \cite{wei2022chain} משפר ביצועים על ידי הרחבת
הטקסט הביניים בחלון ההקשר --- אך מגדיל צריכת טוקנים. איזון בין
איכות לעלות הוא החלטת ארכיטקטורה.
\end{insightbox}

\subsection{השלכות על עלות}

כל טוקן בחלון ההקשר --- קלט ופלט --- עולה כסף במודלים מסחריים.
שליפת מסמכים מיותרים ל-\eng{RAG}, צילומי מסך מדפדפן והיסטוריה ארוכה
מנפחים את העלות. לכן מעבר מ-\eng{GUI} ל-\eng{API}/טקסט הוא גם החלטה
כלכלית.
